\chapter{Công Việc Đầu Tiên Và Cú Sốc Thực Tế}
\markboth{Công Việc Đầu Tiên Và Cú Sốc Thực Tế}{The First Job and the Shock of Reality}

\section{Công Việc Đầu Tiên Thường Không Giống Trong Tưởng Tượng}
\begin{truyen}{Công Việc Đầu Tiên Thường Không Giống Trong Tưởng Tượng}{The First Job Is Usually Not What You Imagined}
\chuhoa{N}{hiều người trẻ đi làm với ý nghĩ rằng mình sẽ được làm những việc quan trọng ngay.}
\textit{(Many young people start work thinking they will handle important tasks right away.)}

Nhưng việc đầu tiên thường là việc nhỏ, lặp lại, và có phần chán.
\textit{(But the first work is often small, repetitive, and a little boring.)}

Đó không hẳn là coi thường em.
\textit{(That is not always disrespect.)}

Đó là cách đời kiểm tra xem em có làm tốt việc nhỏ không.
\textit{(It is how life checks whether you can do small work well.)}

Ai cũng muốn được tin ngay.
\textit{(Everyone wants quick trust.)}

Nhưng phần lớn niềm tin được xây từ những việc ít ai để ý.
\textit{(But most trust is built from tasks that few people notice.)}

\end{truyen}

\begin{baihoc}
Đừng khinh việc nhỏ ở công việc đầu tiên. Rất nhiều người mất cơ hội lớn chỉ vì làm ẩu những việc tưởng như bé.
\textit{(Do not despise small tasks in your first job. Many people lose big chances because they do small work carelessly.)}
\end{baihoc}

\begin{ghinhoanh}
\textbf{Short English to remember:}

Do not despise small tasks in your first job.

Many people lose big chances because they do small work carelessly.

\end{ghinhoanh}

\begin{tuhocnhanh}
\textbf{Cau hoi tu hoc / Self-study questions}

1. Van de chinh la gi? \textit{What is the main problem?}

2. Bai hoc thuc te la gi? \textit{What is the practical lesson?}

3. Mot cau tieng Anh em muon nho la cau nao? \textit{Which English sentence do you want to remember?}
\end{tuhocnhanh}
\ngancach

\section{Đi Làm Mới Biết Thông Minh Không Đủ}
\begin{truyen}{Đi Làm Mới Biết Thông Minh Không Đủ}{Work Teaches You That Intelligence Is Not Enough}
\chuhoa{M}{ột người có thể hiểu rất nhanh nhưng vẫn làm đồng đội mệt nếu trễ giờ, khó hợp tác, và không giữ lời.}
\textit{(A person may understand things fast and still exhaust a team by being late, hard to work with, and careless with promises.)}

Ngoài đời, người ta không chỉ nhìn em biết gì.
\textit{(In real life, people do not only see what you know.)}

Người ta còn nhìn em làm việc cùng người khác thế nào.
\textit{(They also see how you work with others.)}

Có những người rất sáng nhưng đi chậm vì kỹ năng người quá yếu.
\textit{(Some very bright people move slowly because their people skills are weak.)}

Thông minh là lợi thế lớn.
\textit{(Intelligence is a major advantage.)}

Nhưng tính ổn định và đáng tin mới giúp em đi lâu.
\textit{(But steadiness and trustworthiness are what help you last.)}

\end{truyen}

\begin{baihoc}
Muốn đi xa, người trẻ phải rèn cả năng lực lẫn cách làm việc với người khác. Một mình giỏi chưa chắc đã đủ.
\textit{(To go far, young people must build both skill and the ability to work with others. Being clever alone is not always enough.)}
\end{baihoc}

\begin{ghinhoanh}
\textbf{Short English to remember:}

To go far, young people must build both skill and the ability to work with others.

Being clever alone is not always enough.

\end{ghinhoanh}

\begin{tuhocnhanh}
\textbf{Cau hoi tu hoc / Self-study questions}

1. Van de chinh la gi? \textit{What is the main problem?}

2. Bai hoc thuc te la gi? \textit{What is the practical lesson?}

3. Mot cau tieng Anh em muon nho la cau nao? \textit{Which English sentence do you want to remember?}
\end{tuhocnhanh}
\ngancach

\section{Làm Lại Từ Đầu Là Chuyện Bình Thường}
\begin{truyen}{Làm Lại Từ Đầu Là Chuyện Bình Thường}{Doing It Again from the Start Is Normal}
\chuhoa{N}{gày đi học, làm gần đúng đôi khi vẫn được điểm phần nào.}
\textit{(In school, being almost right sometimes still earns some marks.)}

Đi làm rồi, một việc gần đúng có thể vẫn là chưa dùng được.
\textit{(At work, something almost right may still be unusable.)}

Nhiều người trẻ tự ái khi bị yêu cầu làm lại.
\textit{(Many young people feel offended when asked to redo work.)}

Họ tưởng đó là phủ nhận mình.
\textit{(They think it is rejection.)}

Thực ra nhiều lần đó chỉ là tiêu chuẩn.
\textit{(In many cases it is simply the standard.)}

Việc chưa đạt thì phải sửa cho đạt.
\textit{(If the work is not good enough, it must be fixed.)}

\end{truyen}

\begin{baihoc}
Bớt tự ái với phản hồi sẽ giúp em lớn nhanh hơn. Làm lại cho tốt không làm em nhỏ đi.
\textit{(If you become less defensive about feedback, you grow faster. Reworking something well does not make you smaller.)}
\end{baihoc}

\begin{ghinhoanh}
\textbf{Short English to remember:}

If you become less defensive about feedback, you grow faster.

Reworking something well does not make you smaller.

\end{ghinhoanh}

\begin{tuhocnhanh}
\textbf{Cau hoi tu hoc / Self-study questions}

1. Van de chinh la gi? \textit{What is the main problem?}

2. Bai hoc thuc te la gi? \textit{What is the practical lesson?}

3. Mot cau tieng Anh em muon nho la cau nao? \textit{Which English sentence do you want to remember?}
\end{tuhocnhanh}
\ngancach

\section{Sếp Không Phải Lúc Nào Cũng Là Thầy Tử Tế}
\begin{truyen}{Sếp Không Phải Lúc Nào Cũng Là Thầy Tử Tế}{A Boss Is Not Always a Kind Teacher}
\chuhoa{N}{hiều người bước vào đời và sốc vì phát hiện không phải ai hơn mình cũng có trách nhiệm nâng đỡ mình.}
\textit{(Many people enter adult life and are shocked to learn that not everyone above them feels responsible for helping them grow.)}

Có sếp chỉ quan tâm kết quả.
\textit{(Some bosses only care about results.)}

Có người thiếu kiên nhẫn.
\textit{(Some lack patience.)}

Có người giỏi việc nhưng dở cách dẫn người.
\textit{(Some are good at the job but poor at leading people.)}

Điều đó không đẹp, nhưng có thật.
\textit{(That is not beautiful, but it is real.)}

Người trẻ cần học cách tự học cả trong môi trường không lý tưởng.
\textit{(Young people need to learn even inside environments that are not ideal.)}

\end{truyen}

\begin{baihoc}
Ra đời là chấp nhận rằng không phải lúc nào em cũng gặp người hướng dẫn tốt. Càng vậy, năng lực tự học càng quý.
\textit{(Adult life means accepting that you will not always meet good mentors. That is why self-learning becomes more valuable.)}
\end{baihoc}

\begin{ghinhoanh}
\textbf{Short English to remember:}

Adult life means accepting that you will not always meet good mentors.

That is why self-learning becomes more valuable.

\end{ghinhoanh}

\begin{tuhocnhanh}
\textbf{Cau hoi tu hoc / Self-study questions}

1. Van de chinh la gi? \textit{What is the main problem?}

2. Bai hoc thuc te la gi? \textit{What is the practical lesson?}

3. Mot cau tieng Anh em muon nho la cau nao? \textit{Which English sentence do you want to remember?}
\end{tuhocnhanh}
\ngancach

\section{Đi Làm Mới Biết Mệt Kiểu Khác Đi Học}
\begin{truyen}{Đi Làm Mới Biết Mệt Kiểu Khác Đi Học}{Work Exhausts You in a Different Way from School}
\chuhoa{Đ}{i học mệt vì kiểm tra, vì bài vở, vì áp lực thành tích.}
\textit{(School is tiring because of tests, lessons, and performance pressure.)}

Đi làm mệt thêm vì lặp lại, vì trách nhiệm, vì phải giữ thái độ kể cả lúc không vui.
\textit{(Work adds new fatigue: repetition, responsibility, and the need to stay professional even when unhappy.)}

Có người không chuẩn bị tâm lý cho kiểu mệt này nên rất dễ vỡ mộng.
\textit{(Some people do not prepare mentally for this kind of tiredness, so they become disillusioned very quickly.)}

Mỗi giai đoạn của đời có một kiểu nặng riêng.
\textit{(Every stage of life has its own weight.)}

Khôn ngoan là học cách mang nó đúng, không phải ngây thơ nghĩ nó sẽ không tới.
\textit{(Wisdom means learning how to carry it well, not pretending it will never come.)}

\end{truyen}

\begin{baihoc}
Đi làm không chỉ cần năng lực. Nó còn cần sức bền tâm lý cho những ngày rất bình thường nhưng rất hao người.
\textit{(Work requires more than ability. It also needs mental stamina for ordinary days that quietly drain you.)}
\end{baihoc}

\begin{ghinhoanh}
\textbf{Short English to remember:}

Work requires more than ability.

It also needs mental stamina for ordinary days that quietly drain you.

\end{ghinhoanh}

\begin{tuhocnhanh}
\textbf{Cau hoi tu hoc / Self-study questions}

1. Van de chinh la gi? \textit{What is the main problem?}

2. Bai hoc thuc te la gi? \textit{What is the practical lesson?}

3. Mot cau tieng Anh em muon nho la cau nao? \textit{Which English sentence do you want to remember?}
\end{tuhocnhanh}
\ngancach

\section{Văn Phòng Cũng Có Chính Trị}
\begin{truyen}{Văn Phòng Cũng Có Chính Trị}{Even Offices Have Politics}
\chuhoa{Đ}{i làm một thời gian, người trẻ mới thấy năng lực không phải thứ duy nhất vận hành nơi làm việc.}
\textit{(After some time at work, young people see that skill is not the only force in a workplace.)}

Họ tưởng cứ làm tốt là đủ để mọi việc chạy thẳng và sạch.
\textit{(They think doing good work is enough for everything to stay straight and fair.)}

Nhưng còn có quan hệ, cách nói, thời điểm, và cách người ta nhìn nhau.
\textit{(But there are also relationships, timing, communication, and perception.)}

Hiểu điều đó không phải để gian hơn, mà để bớt ngây thơ hơn.
\textit{(Understanding that is not for becoming cynical, but for becoming less naive.)}

\end{truyen}

\begin{baihoc}
Đi làm không chỉ học việc. Nó còn là học cách nhìn một hệ thống thật.
\textit{(Working life teaches not only tasks, but also how a real system behaves.)}
\end{baihoc}

\begin{ghinhoanh}
\textbf{Short English to remember:}

Working life teaches not only tasks, but also how a real system behaves.

\end{ghinhoanh}

\begin{tuhocnhanh}
\textbf{Cau hoi tu hoc / Self-study questions}

1. Van de chinh la gi? \textit{What is the main problem?}

2. Bai hoc thuc te la gi? \textit{What is the practical lesson?}

3. Mot cau tieng Anh em muon nho la cau nao? \textit{Which English sentence do you want to remember?}
\end{tuhocnhanh}
\ngancach

\section{Hỏi Không Phải Là Kém}
\begin{truyen}{Hỏi Không Phải Là Kém}{Asking Questions Is Not Weakness}
\chuhoa{N}{hiều nhân viên mới sợ hỏi vì ngại bị nghĩ là chậm hay dốt.}
\textit{(Many new workers fear asking questions because they do not want to look slow or foolish.)}

Họ tưởng im lặng sẽ giúp mình giữ hình ảnh thông minh.
\textit{(They think silence protects an image of intelligence.)}

Nhưng im lặng sai chỗ thường làm lỗi lớn hơn và mất thời gian hơn.
\textit{(But silence in the wrong place usually creates larger mistakes and wastes more time.)}

Người học nhanh không phải người giấu việc không hiểu.
\textit{(Fast learners are not those who hide confusion.)}

Họ là người hỏi đúng lúc rồi sửa nhanh.
\textit{(They ask at the right time and correct quickly.)}

\end{truyen}

\begin{baihoc}
Hỏi sớm thường cứu được cả việc lẫn lòng tự trọng về sau.
\textit{(Asking early often saves both the work and your dignity later.)}
\end{baihoc}

\begin{ghinhoanh}
\textbf{Short English to remember:}

Asking early often saves both the work and your dignity later.

\end{ghinhoanh}

\begin{tuhocnhanh}
\textbf{Cau hoi tu hoc / Self-study questions}

1. Van de chinh la gi? \textit{What is the main problem?}

2. Bai hoc thuc te la gi? \textit{What is the practical lesson?}

3. Mot cau tieng Anh em muon nho la cau nao? \textit{Which English sentence do you want to remember?}
\end{tuhocnhanh}
\ngancach

\section{Đi Trễ Một Lần Cũng Dạy Rất Nhiều}
\begin{truyen}{Đi Trễ Một Lần Cũng Dạy Rất Nhiều}{One Late Arrival Can Teach a Lot}
\chuhoa{C}{ó người chỉ đến trễ một lần trong buổi họp đầu nhưng nhớ mãi cảm giác ánh mắt mọi người nhìn mình.}
\textit{(Someone arrives late once to an early meeting and never forgets the room looking at them.)}

Họ từng tưởng trễ vài phút chỉ là chuyện nhỏ.
\textit{(They once thought a few late minutes were a small issue.)}

Ngoài đời, đúng giờ là tín hiệu của sự tôn trọng trước khi ai đó kịp thấy năng lực của em.
\textit{(In adult life, punctuality signals respect before anyone can judge your ability.)}

Một lỗi nhỏ về giờ giấc đôi khi làm người khác đọc ra cả một thái độ.
\textit{(A small time error can make others read a larger attitude.)}

\end{truyen}

\begin{baihoc}
Đi làm dạy rằng nhiều điều bé lại mang ý nghĩa lớn hơn ta tưởng.
\textit{(Work teaches that small things often carry larger meaning than we expect.)}
\end{baihoc}

\begin{ghinhoanh}
\textbf{Short English to remember:}

Work teaches that small things often carry larger meaning than we expect.

\end{ghinhoanh}

\begin{tuhocnhanh}
\textbf{Cau hoi tu hoc / Self-study questions}

1. Van de chinh la gi? \textit{What is the main problem?}

2. Bai hoc thuc te la gi? \textit{What is the practical lesson?}

3. Mot cau tieng Anh em muon nho la cau nao? \textit{Which English sentence do you want to remember?}
\end{tuhocnhanh}
\ngancach

\section{Ngày Bình Thường Mới Xây Uy Tín}
\begin{truyen}{Ngày Bình Thường Mới Xây Uy Tín}{Ordinary Days Build Trust}
\chuhoa{K}{hông phải ngày bùng nổ nào cũng quyết định chỗ đứng của em trong công việc đầu tiên.}
\textit{(Not every dramatic day decides your place in a first job.)}

Nhiều người trẻ tưởng phải làm một việc thật lớn mới được tin.
\textit{(Many young people think they must do something huge to earn trust.)}

Nhưng niềm tin thường được ghép từ những ngày rất bình thường em vẫn làm việc gọn, đúng, và đều.
\textit{(But trust is usually built from ordinary days when you work neatly, correctly, and steadily.)}

Người đi được lâu là người giữ chất lượng cả khi chẳng ai chú ý.
\textit{(Those who last are those who keep quality even when nobody is watching.)}

\end{truyen}

\begin{baihoc}
Việc đầu tiên không chỉ kiểm tra tài. Nó kiểm tra độ đều của em.
\textit{(The first job does not only test talent. It tests consistency.)}
\end{baihoc}

\begin{ghinhoanh}
\textbf{Short English to remember:}

The first job does not only test talent.

It tests consistency.

\end{ghinhoanh}

\begin{tuhocnhanh}
\textbf{Cau hoi tu hoc / Self-study questions}

1. Van de chinh la gi? \textit{What is the main problem?}

2. Bai hoc thuc te la gi? \textit{What is the practical lesson?}

3. Mot cau tieng Anh em muon nho la cau nao? \textit{Which English sentence do you want to remember?}
\end{tuhocnhanh}
\ngancach

\section{Nghỉ Việc Hay Ở Lại}
\begin{truyen}{Nghỉ Việc Hay Ở Lại}{Should You Quit or Stay?}
\chuhoa{C}{ó người đi làm vài tháng đã thấy nản và nghĩ đến chuyện bỏ ngay.}
\textit{(Some people feel discouraged after a few months of work and immediately think about quitting.)}

Họ tưởng cứ mệt là chắc chắn mình chọn sai.
\textit{(They think feeling tired automatically means the choice was wrong.)}

Nhưng có lúc mệt vì đang lớn, có lúc mệt vì môi trường thực sự không đáng ở.
\textit{(But sometimes you are tired because you are growing, and sometimes because the environment is truly unhealthy.)}

Khó nhất là học cách phân biệt hai kiểu mệt đó.
\textit{(The hardest part is learning to tell those two kinds of tiredness apart.)}

\end{truyen}

\begin{baihoc}
Người trưởng thành không ở lì cũng không bỏ chạy quá sớm. Họ nhìn cho rõ rồi mới quyết.
\textit{(Mature people neither stay blindly nor run too early. They look carefully before deciding.)}
\end{baihoc}

\begin{ghinhoanh}
\textbf{Short English to remember:}

Mature people neither stay blindly nor run too early.

They look carefully before deciding.

\end{ghinhoanh}

\begin{tuhocnhanh}
\textbf{Cau hoi tu hoc / Self-study questions}

1. Van de chinh la gi? \textit{What is the main problem?}

2. Bai hoc thuc te la gi? \textit{What is the practical lesson?}

3. Mot cau tieng Anh em muon nho la cau nao? \textit{Which English sentence do you want to remember?}
\end{tuhocnhanh}
